\section{研究课题识别}

\subsection{研究背景与动机}

上文对文献的梳理揭示了位于DeFi借贷风险、行为金融和链上信用评估三者交叉处的一个明确的尚未充分探索的研究问题。本节明确本研究拟解决的具体研究课题，包括其范围、研究问题和可检验假设，以指导实证检验。

本研究的提出基于如下观察：DeFi借贷协议生成了每位参与者每一次操作的综合且带有时间戳的记录，然而现有文献主流将这些数据视为截面风险分类的静态特征来源，而不是从数据中窥探借款人在管理（或未能管理）不断累积的风险暴露中经历过的动态过程的窗口。这无疑是一次错失良机，因为借款人行为的时效维度——包括主动调整的顺序、时间和幅度相对于头寸健康状况的变化——很可能包含信用相关信息，而这些信息是抵押率、健康因子或其他传统风险指标的时间点快照所无法捕捉的。

因此，研究课题定义如下：\textbf{DeFi借贷协议中，借款人对头寸风险不断上升所作的主动行为调整与后续清算事件或其他不利风险结果之间的经验关系。}该课题位于金融风险评估的更广泛领域之中，但其通过对行为过程变量（而非状态变量）的强调、对公开可核实链上数据的利用以及对行为金融理论框架的基础性依归，使自身于此区别分析。

\subsection{研究问题}

本研究解决递进层级结构的两个研究问题：

\begin{description}[leftmargin=*,labelindent=\parindent, itemindent=0pt]
\item[RQ1（行为层面）。] 在DeFi借款人的贷款头寸接近协议定义的清算阈值时，是否会表现出系统性非随机的主动调整模式——包括追加抵押品、偿还债务、平仓或继续承担风险——这些模式能否与价格变动的机械效应和清算机器人操作形成区分？

\item[RQ2（信用层面）。] 在控制常规链上风险指标（包括当前健康因子、抵押率、账户活动、资产集中度和市场波动性）的条件下，捕捉主动调整\emph{轨迹}的行为过程变量是否为未来清算事件提供了统计和经济上显著的增量预测能力？
\end{description}

RQ1询问是否存在一个值得解释的行为现象：这是一个关于存在、量级和主动借贷者对风险积累响应的结构的问题。即使缺少信用预测，该问题也具有固有的描述性价值。RQ2则询问这一现象是否对风险评估具有实际意义：即行为信息是否能在直接、简单和现成的风险指标之外提供超越值。研究问题的层级结构意味着RQ2以RQ1为条件：如果不存在系统性的行为响应，那么这些响应能否预测结果的问题便无意义。

\subsection{研究假说}

研究问题可转化为如下可检验假说：

\subsubsection*{RQ1：行为假说}

\begin{description}[leftmargin=*,labelindent=\parindent, itemindent=0pt]
\item[H1a（主动调整的发生概率）。] 在控制标的抵押资产同一时期价格变动的条件下，观察到借款人追加抵押品或偿还债务的概率随健康因子由上而下接近清算阈值而单调递增。

\item[H1b（损失厌恶非对称性）。] 在健康因子每单位恶化的情况下（损失域），借款人主动风险减轻行为（追加抵押品、偿还债务）的绝对量级，大于在健康因子每单位改善的情况下（收益域）追加借款或提取抵押品的绝对量级，符合前景理论的损失厌恶预测。

\item[H1c（参照点断点）。] 在对应清算阈值的健康因子值处存在主动风险减轻行为发生概率的断点，该断点显著超越于仅使用抵押品价格变动所能预测的平滑关系。
\end{description}

\subsubsection*{RQ2：信用预测假说}

\begin{description}[leftmargin=*,labelindent=\parindent, itemindent=0pt]
\item[H2a（增量预测能力）。] 在预测清算事件的样本外预测能力方面，包含行为过程变量（特别是头寸接近清算阈值时的主动调整的存在、时机和量级）的模型实现了优于仅使用常规链上风险指标的基准模型的统计显著提升。

\item[H2b（时序优先性）。] 在时期\(t-1\)至\(t\)度量的行为调整变量，即使在控制截至时间\(t\)的健康因子及其近期变化轨迹的条件下，也能改善对时期\(t\)至\(t+k\)清算事件的预测能力。
\end{description}

\subsection{研究定位与创新性}

所提出的研究在三个方面对现有文献做出清晰贡献。第一，它将DeFi借贷研究的分析着力点从协议层面风险转向借款人层面行为，为理解DeFi市场参与者如何响应智能合约参数生成的激励结构提供了微观基础。第二，它通过在一个自然发生的金融情境中检验前景理论的预测——其中清算阈值这一参照点是客观、可观察且对所有参与者均相同的——连接了DeFi风险建模和行为金融这两个基本独立的文献群体，从而应对了限制许多行为金融研究外部有效性的核心度量挑战~\cite{barberis2013thirty}。第三，它引入了一种过程视角的信用风险评估方法，补充了迄今为止主导链上信用风险文献的静态性和基于特征的方法~\cite{ghosh2024onchain}，可能为DeFi借贷协议构建行为性早期预警指标开辟新途径。

有必要清晰的界定本研究\emph{不}主张什么。它不声称DeFi借款人是非理性的，也不声称他们的行为必然偏离理性预期模型的预测。经验方法的设计目的是检验行为变量是否在常规风险指标之外\emph{增加了信息}，而非在不同理性模型之间加以裁决。研究假说明确地以\emph{增量}预测能力为框架，而非绝对准确度。本研究也不声称任何观察到的行为模式必然可以推广至其他DeFi协议、时间段或市场条件；外部有效性是后续重复和扩展的任务。

%
% Figure 3: Risk Trajectory / Health Factor Illustration
%
\begin{figure}[t]
\centering
\resizebox{\textwidth}{!}{%
\begin{tikzpicture}[
  scale=1.0,
  >=stealth,
  axisline/.style={->, line width=0.6pt},
  dataline/.style={line width=0.8pt},
]
  % Axes
  \draw[axisline] (0,0) -- (9.5,0) node[right, font=\footnotesize] {Time $\longrightarrow$};
  \draw[axisline] (0,0) -- (0,5.2) node[above, font=\footnotesize] {Health Factor};
  
  % Labels
  \node[font=\scriptsize] at (4.5,-0.5) {Observation Window (daily frequency)};
  
  % Y-axis labels
  \draw[dashed, gray!40] (0,0.0) -- (8.5,0.0);
  \node[font=\scriptsize, left] at (-0.15,0.0) {0};
  \draw[dashed, gray!40] (0,1.0) -- (8.5,1.0);
  \node[font=\scriptsize, left] at (0,1.0) {1.0};
  \node[font=\tiny, red!70!black] at (-0.7, 0.55) {\shortstack{Liq.\\Thresh.}};
  
  \draw[dashed, gray!40] (0,2.0) -- (8.5,2.0);
  \node[font=\scriptsize, left] at (0,2.0) {2.0};
  \draw[dashed, gray!40] (0,3.0) -- (8.5,3.0);
  \node[font=\scriptsize, left] at (0,3.0) {3.0};
  \draw[dashed, gray!40] (0,4.0) -- (8.5,4.0);
  \node[font=\scriptsize, left] at (0,4.0) {4.0};
  
  % Risk bands (shaded)
  \fill[red!8] (0,1.0) rectangle (8.5,0.0);
  \fill[orange!8] (0,2.0) rectangle (8.5,1.0);
  \fill[yellow!8] (0,3.0) rectangle (8.5,2.0);
  \fill[green!8] (0,5.0) rectangle (8.5,3.0);
  
  \node[font=\tiny\sffamily, text=red!50!black, opacity=0.6] at (1.2,0.5) {Liquidation Zone};
  \node[font=\tiny\sffamily, text=orange!50!black, opacity=0.6] at (1.2,1.5) {High Risk (HF $<$ 2)};
  \node[font=\tiny\sffamily, text=yellow!50!black, opacity=0.6] at (1.2,2.5) {Moderate Risk};
  \node[font=\tiny\sffamily, text=green!50!black, opacity=0.6] at (1.2,4.0) {Safe Zone (HF $>$ 3)};
  
  % Simulated HF trajectory
  \draw[dataline, blue!60, line width=1.2pt] plot[smooth] coordinates {
    (0,3.5) (0.35,3.3) (0.7,3.1) (1.0,2.8) (1.4,2.5) (1.7,2.3) 
    (2.1,2.0) (2.3,1.7) (2.4,1.5) (2.55,1.25) 
  };
  
  % Active adjustment events (stars)
  \node[star, star points=5, star point ratio=2.25, draw=green!70!black, fill=green!60, minimum size=2mm, inner sep=0pt] at (1.4,2.5) {};
  \draw[->, green!60!black, line width=0.5pt] (1.4,2.9) -- (1.4,2.5) node[pos=0, above, font=\tiny] {+collateral};
  
  \node[star, star points=5, star point ratio=2.25, draw=green!70!black, fill=green!60, minimum size=2mm, inner sep=0pt] at (2.1,2.0) {};
  \draw[->, green!60!black, line width=0.5pt] (2.1,2.4) -- (2.1,2.0) node[pos=0, above, font=\tiny] {repay debt};
  
  % Liquidation event
  \node[circle, draw=red!80!black, fill=red!60, minimum size=2.5mm, inner sep=0pt] at (2.55,1.25) {\tiny\color{white}\textbf{!}};
  \draw[->, red!80!black, line width=0.5pt] (2.9,1.25) -- (2.65,1.25) node[pos=0, right, font=\tiny] {liquidation};
  
  % Legend (inside plot area, top right)
  \node[font=\tiny, draw=black!30, fill=white, rounded corners=2pt, inner sep=3pt] at (6.8,4.5) {
    \shortstack[l]{
      \tikz{\draw[blue!60, line width=1.1pt] (0,0)--(0.4,0);} \, HF trajectory\\
      \tikz{\node[star, star points=5, star point ratio=2.25, fill=green!60, minimum size=1.3mm, inner sep=0pt] at (0,0) {};} \, Active adjustment\\
      \tikz{\node[circle, fill=red!60, minimum size=1.3mm, inner sep=0pt] at (0,0) {};} \, Liquidation event
    }
  };

\end{tikzpicture}%
}
\caption{Illustrative health factor (HF) trajectory and behavioral response. A borrowing position's HF declines toward the liquidation threshold (HF = 1.0) as collateral value decreases. Active borrower adjustments (green stars) represent collateral additions or debt repayment. The central empirical question is whether the pattern of these adjustments—their timing, magnitude, and frequency—provides predictive information about eventual liquidation beyond what is captured by the HF itself.}
\label{fig:trajectory}
\end{figure}

\vspace{4pt}
